\chapter{Data Engineering for Pretraining}
\label{ch:data-engineering}

% Part II, Chapter 9: Filtering, deduplication, mixing, packing, and
% contamination.
%
% Core question: Chapter 8 told you that you need 6B tokens. Where do they
% come from, and how do you make sure they're worth learning from?
%
% Planned sections:
%   9.1  Hook: "The Model Can Only Learn What the Data Contains."
%   9.2  Web Crawl Pipelines: Common Crawl to Clean Text
%   9.3  Quality Filtering: Perplexity, Classifiers, and Heuristics
%   9.4  Deduplication: Exact (MinHash) and Fuzzy
%   9.5  Domain Mixing: Code, Math, Multilingual, and Proportions
%   9.6  Document Packing and Attention Masks
%   9.7  Data Curriculum: Ordering, Staging, and Annealing
%   9.8  Data Contamination: When Training Data Leaks into Benchmarks
%   9.9  Failure Modes in Data Engineering
%
% Key thesis: LLaMA's core contribution was not architecture innovation
% (there was almost none). It was data engineering. At the same parameter
% count, better data = better model. This is the most underrated lesson
% of 2023--2024.
%
% Lab 9 (planned): Take a raw Common Crawl shard, apply filtering and dedup,
% measure how much data survives, train two small models (raw vs cleaned),
% compare quality.

\textit{This chapter is a placeholder. It will cover web crawl pipelines,
quality filtering, deduplication, domain mixing, document packing, data
curriculum, and data contamination.}
